\documentclass[11pt]{article}
\usepackage[margin=1in]{geometry}
\usepackage{amsmath,amssymb}
\usepackage[hidelinks]{hyperref}
\usepackage[round]{natbib}

\title{Asymmetric Capital Control: Executive Summary}
\date{}

\begin{document}
\maketitle
\thispagestyle{empty}

\noindent
Codes in parentheses refer to results in the proofs document, which
carries an index of them. Claims not so marked are either framing or are
attributed to the literature.

\bigskip

Two banks facing identical regulation can run visibly different capital
policies, and the difference matters: one that treats a shortfall against
its own target as far worse than an equivalent \textbf{surplus} is good will hold
more capital, pay out later, and recapitalise sooner than the rules alone
require. Whether that conservatism comes from the regulation or from the
bank's own asymmetric attitude to shortfalls determines how it responds
when the rules change---and the two are not distinguishable by inspection,
because both produce a bank holding more capital than the minimum.
Recovering the preference separately from the regulation is therefore the
prior question, and this paper's central finding is that the asymmetric
aversion is invisible in the natural \textbf{volatility-based measure} of it, and
visible instead in the thresholds the bank commits to. In a solved model
of a bank managing its \textbf{capital buffer} under a \textbf{regulatory cap}, the
\textbf{definitional evaluation} of the \textbf{deficit}-to-\textbf{surplus} \textbf{volatility ratio},
computed directly from the model's structure rather than estimated,
depends only on \textbf{cap geometry}---how the \textbf{regulatory cap} varies with the
\textbf{buffer}---and carries no information about the preference at all~(SCG). The
preference does show up in where the bank commits to act, and both
responses are obtained in closed form: a more shortfall-averse bank
retains a larger \textbf{buffer} before paying out at the \textbf{dividend barrier},
and recapitalises earlier, from a higher level of capital, at the
\textbf{recapitalisation trigger}~(TCS). The \textbf{fixed issuance cost} moves that
trigger the opposite way~(KCS), and the opposition of the two is
decisive: it renders the Jacobian of the two thresholds in the two
parameters non-singular, so preference and cost are jointly identified
from the boundaries alone~(IDN). Much of empirical finance reads a firm's
risk attitude off its volatility; in this model that \textbf{identification}
fails outright, and the decision boundaries carry the information
instead.

Two structural facts underpin that finding. First, setting the asymmetry
parameter to its neutral value removes the \textbf{shortfall penalty} entirely,
giving an \textbf{exact nesting} of the \textbf{risk-neutral benchmark} rather than a
limiting approximation of it~(NST), so the parameter measures asymmetry in
excess of what regulation and costs alone would produce. Second, because
the \textbf{shortfall penalty} is kinked-linear rather than S-shaped, the running
payoff stays concave for every degree of asymmetry~(RPC). This is the
paper's departure from \textbf{reference}-dependent models built on the S-shaped
\textbf{value function} of prospect theory \citep{kahneman1979}, which are convex
below the \textbf{reference} and therefore require a concavification argument
before any threshold structure can be recovered at all. The claim is
stated at exactly that strength and no further: it concerns the
objective's flow term, and does not deliver a concave \textbf{value function}. The
\textbf{value function} is in fact locally convex above the \textbf{recapitalisation}
trigger~(LNC), but that non-concavity is caused by the \textbf{fixed issuance cost}
and is present even at the neutral parameter value, so it is localised and
cost-driven rather than global and preference-driven---precisely the
difference from the S-shaped case.

A \textbf{discrete-time model} is worked out first and closed completely, with the
\textbf{capital buffer} written as a ratio of equity to \textbf{required capital}. The
\textbf{continuous-time formulation} restates the same bank's problem with three
controls of distinct type: the \textbf{classical control} of \textbf{risky exposure},
constrained by the \textbf{regulatory cap}; the \textbf{singular control} of \textbf{dividends},
acting at the \textbf{dividend barrier}; and the \textbf{impulse control} of
\textbf{recapitalisation}, carrying both a fixed and a \textbf{proportional issuance cost}.

The \textbf{regulatory cap} is not a single constraint but two regimes, a solvency
branch and a liquidity branch, with the binding one switching at the level
where they cross. The bank's \textbf{capped exposure} is the smaller of the two at
each point, which gives the cap a kink of its own, distinct from the kink
in the preference. A consistency condition on the parameters governs
behaviour at the \textbf{distress boundary}: the drift there is positive only if
the return on assets exceeds the return on liabilities, which the
parameters used throughout satisfy~(CCP). This is what pushes the \textbf{buffer}
away from distress rather than into it, and it is the reason the distress
boundary is never reached under the solved policy.

An \textbf{accounting identity} ties the bank's \textbf{return on equity} to its growth and
retention~(ACI). \textbf{Deficits} are taken to decay more slowly than \textbf{surpluses}, a
\textbf{persistence asymmetry}~(PSA), and from that the central discrete-time
result follows: the \textbf{volatility ratio} is a fixed power of the asymmetry
parameter~(FPV), depending only on the \textbf{saturation limits} of the transition
between the two states and not on how the transition is shaped in
between~(ROB). The same recursion is observationally equivalent to a
standard \textbf{asymmetric volatility specification} of the EGARCH class~(EGE), so
the ratio could in principle be recovered from returns without observing
\textbf{required capital} directly, and a first-passage argument gives the sojourn
time---the expected length of time the \textbf{buffer} spends below the
\textbf{reference}---in closed form~(SJT).

In continuous time, a \textbf{liquidation payoff} bounds the \textbf{value function} from
below~(LQB), and a change of variable shows that the discrete-time state
and the continuous-time \textbf{capital buffer} are one state measured from one
origin~(TSO). The apparent breakdown of the model at the distress
boundary, the point at which the \textbf{buffer} exactly matches \textbf{required capital},
turns out to be a \textbf{coordinate artefact} rather than a feature of the
problem: in the reconciled coordinate the model is well-behaved
everywhere, and the boundary is never actually reached~(CDA). That
reconciliation is what makes it meaningful to compare the same
preference's legibility in the thresholds and in realised volatility on
equal footing.

A separate finding concerns where the \textbf{regulatory cap} actually binds. Under
the solved policy the cap is active across the whole \textbf{inaction region}: the
solver searches the exposure over a grid spanning the admissible range and
retains the maximiser, so a slack optimum would have been found had one
existed, and at almost every solved case none is~(CBI). What closes the
problem is therefore the \textbf{impulse control} rather than any interior
relaxation of the cap. The paper states this at the strength it has---a
finding at these parameters rather than a theorem, with the formulation
accommodating either outcome without modification---and the solver
producing it is separately documented and validated~(SLV).

The thresholds' \textbf{comparative statics} give the fuller picture, and they
are theorems rather than observations. Because the asymmetry parameter
enters the objective linearly and never restricts the admissible set, the
sensitivity of value to it is the expected discounted shortfall incurred
under the optimal policy, and every boundary response reduces to a
property of that single object~(ENV). Both smooth-fit conditions then
supply its boundary data, and the derivatives follow in closed form:
$dy^*/d\lambda_S$ from the second-order fit at the barrier ---
the first-order condition being degenerate there --- and $dx_L/d\lambda_S$
from a first-order expansion at the hitting time~(TCS). The
\textbf{fixed issuance cost} admits the identical treatment with the discounted
count of injections in place of the shortfall, and the signs oppose on the
trigger: a steeper penalty makes the bank intervene earlier, a larger
fixed cost makes it wait longer, while both raise the barrier~(KCS). That
opposition is what makes the two parameters separately recoverable, and it
does so unconditionally rather than by comparing magnitudes~(IDN).
Empirically the two write themselves into different observables ---
preference moves the levels, the \textbf{fixed issuance cost} moves the
\textbf{trigger-to-target gap}, which vanishes exactly at zero fixed cost~(KSW). That threshold geometry should carry preference
information is not without precedent at the firm level---\citet{angelis2025}
derives an inaction band with two direction-dependent targets from loss
aversion in a price-setting problem---but that model is a discrete
behavioural one rather than a continuous-time control problem, and does
not run the inverse mapping from band to preference.

What the results rest on is stated rather than assumed. Differentiability
of the value in each parameter needs no hypothesis at all: the value is a
supremum of functions affine in that parameter, hence convex in it, hence
differentiable off a countable set~(REG). Differentiability of the
boundaries reduces to the implicit function theorem, whose non-degeneracy
conditions turn out to be the same ones that sign the derivatives, so each
does double duty~(BDR). The two second-order conditions the identification
needs are likewise theorems: because the impulse operator is affine, the
\textbf{injection target} is independent of where the injection is made from, and
the required curvature at each boundary follows from a sign-change
argument and from the variational inequality on the intervention
region~(SOC). Verification itself is now unconditional: it rests on four
hypotheses, all of which hold, and concavity is not among them --- its
absence is therefore no obstruction~(VER). The one point that had
remained --- whether the \textbf{value function} is continuously
differentiable at the \textbf{recapitalisation trigger} --- is settled by
the same variational inequality: a corner there would have to be convex,
and a convex corner is incompatible with the value function being a
viscosity supersolution, so smooth fit holds and the trigger derivative
is continuous~(SMF). What the numerics could not resolve, because the
discrete scheme carries a corner at the boundary by construction, the
argument settles by showing a corner cannot exist. Each of the
document's thirty-five numbered results carries one canonical verifier,
recorded in its verification appendix: five are proved in Lean with no
unproved gaps, four more have their algebraic cores machine-checked in
Lean with the analytic input stated as an explicit hypothesis, twelve
are exact symbolic checks, and the remainder are proofs in the text or
labelled numerical findings. Separately, whether the
\textbf{induced variance asymmetry}, as opposed to its \textbf{definitional evaluation},
moves with the \textbf{asymmetry parameter} is weakly identified rather than
absent, and the paper reports the magnitude rather than claiming
invariance.

Beyond those, the direction the finding opens is one of scope. The
mechanism rests on a specific choice: the \textbf{kinked-linear penalty}, which
keeps the \textbf{running payoff} concave~(RPC) and, by the same feature, keeps the
thresholds responsive to the preference. Whether committed thresholds
remain informative under other classes of \textbf{shortfall penalty}---and which
classes preserve the contrast between uninformative variances and
informative boundaries---would establish how far the result extends past
the case solved here. \citet{armstrong2019} suggests the question has a
sharp answer: their irrelevance results for constraints require a
disutility that is sublinear in the left tail, a hypothesis a
\textbf{kinked-linear penalty} does not satisfy, which is consistent with limits
binding, and therefore revealing, in exactly this case. Establishing that
connection is future work and not a result of this paper. To our knowledge
no existing firm-level model combines singular payout control with impulse
\textbf{recapitalisation} under an asymmetric preference rather than a constraint;
the recent literature generalising preferences in dividend control remains
singular-only \citep{chen2026,cao2026}. What this paper settles is which
quantities in a solved bank's policy carry the preference and which
cannot, and the answer holds whether or not any bank is ever measured.

\bibliographystyle{plainnat}
\bibliography{references}

\end{document}
